% УВОД, без номер. Ориентировъчен обем: 1 до 2 стр.
% Резюме на дейностите по проекта: защо, каква е целта, какви са задачите,
% какъв е обхватът и как е подредено изложението.
\section*{УВОД}
\label{sec:introduction}
\addcontentsline{toc}{section}{УВОД}

Семантичният уеб описва информацията във вид, който машина може да свърже, вместо
само да покаже. Основната единица на този запис е \term{триплетът}: подлог, сказуемо
и допълнение, всяко от които е ресурс с идентификатор или литерална стойност.
Моделът RDF задава абстрактния синтаксис на този запис~\cite{w3c_rdf11_concepts}, а
езикът SPARQL задава как над множество триплети се формулира заявка~\cite{w3c_sparql11_query}.
Двете спецификации са нормативни и публични, но между тях и работеща програма стои
инженерна задача, която спецификациите съзнателно не решават: как триплетите се
съхраняват, индексират и обхождат така, че заявката да завърши за приемливо време.

Съществуващите реализации решават тази задача, но повечето от тях са сървъри. Те
предполагат отделен процес, мрежов протокол и администриране. Има случаи, в които
това е излишно: настолно приложение, вграден инструмент за анализ, конвейер за
обработка на данни, който трябва да зададе няколко заявки над локален набор от
триплети и да продължи. За такъв случай е подходящо \term{вградено} хранилище, тоест
библиотека, която се свързва към приложението и работи в неговия процес и адресно
пространство.

\textbf{Целта} на курсовия проект е да се проектира и реализира вградено хранилище
за RDF триплети с механизъм за изпълнение на заявки на SPARQL 1.1, и да се измери
поведението му спрямо утвърдени реализации върху едни и същи данни и заявки.

За постигането на целта са формулирани следните \textbf{задачи}:

\begin{enumerate}[topsep=0pt,itemsep=0pt,leftmargin=*]
\item преглед на информационния модел на семантичния уеб, на синтаксисите за обмен
      на RDF и на семантиката на SPARQL, заедно с публикуваните подходи за
      индексиране на триплети;
\item анализ на възможните проектни решения по всяко от отворените места, тоест
      начин на кодиране на термините, набор от индекси, вид на синтактичния
      анализатор и стратегия на планировчика, с обосновка на направения избор;
\item проектиране на слоеста архитектура, в която синтаксисът, съхранението,
      анализът на заявки, планирането и изпълнението са отделени един от друг;
\item програмна реализация на архитектурата на C++20 с изграждане чрез CMake;
\item изграждане на методика за проверка на коректността спрямо публичните тестови
      набори на W3C и за измерване на времето за изпълнение спрямо еталонни
      реализации;
\item провеждане на измерванията, анализ на получените резултати и извеждане на
      заключения за приложимостта на избраните решения.
\end{enumerate}

\textbf{Обхватът} на проекта покрива четене и записване на Turtle и N-Triples,
съхранение с речниково кодиране и пермутирани индекси, синтактичен анализ на
SPARQL 1.1 до алгебрично дърво, планиране на реда на съединенията и изпълнение над
индексите. Правилата за извеждане по RDFS се разглеждат като отделен, незадължителен
слой. Извън обхвата остават протоколът SPARQL през HTTP, езикът SPARQL Update,
разпределеното изпълнение и пълното разсъждаване по OWL 2; последното се разглежда
само описателно, доколкото дисциплината го включва.

\textbf{Предходна работа на автора.} Дипломната работа на автора за ОКС „Бакалавър“
е библиотека на C++ за обектно релационно съответствие, изградена върху изчисления
по време на компилация. Общото с настоящия проект е проектирането на език за заявки
и превръщането му в изпълним план, което тук се пренася върху друг модел на данните.
Изходният текст на предходната работа не се използва.
\TODO{точно заглавие и година на дипломната работа за ОКС „Бакалавър“, за цитиране}

\textbf{Структура на изложението.} Раздел~\ref{sec:literature} разглежда предметната
област и литературата. Раздел~\ref{sec:alternatives} излага разгледаните алтернативи
и обосновава избора. Раздел~\ref{sec:design} представя проектирането, а
Раздел~\ref{sec:implementation} програмната реализация. Раздел~\ref{sec:method}
описва методиката за проверка и измерване, Раздел~\ref{sec:results} получените
резултати.
